\let\negmedspace\undefined
\let\negthickspace\undefined
\documentclass[journal]{IEEEtran}
\usepackage[a5paper, margin=10mm, onecolumn]{geometry}
\usepackage{tfrupee}

\setlength{\headheight}{1cm}
\setlength{\headsep}{0mm}

\usepackage{gvv-book}
\usepackage{gvv}
\usepackage{cite}
\usepackage{amsmath,amssymb,amsfonts,amsthm}
\usepackage{algorithmic}
\usepackage{graphicx}
\usepackage{textcomp}
\usepackage{xcolor}
\usepackage{txfonts}
\usepackage{listings}
\usepackage{enumitem}
\usepackage{mathtools}
\usepackage{gensymb}
\usepackage{comment}
\usepackage[breaklinks=true]{hyperref}
\usepackage{tkz-euclide} 
\usepackage{listings}
\def\inputGnumericTable{}
\usepackage[latin1]{inputenc}
\usepackage{color}
\usepackage{array}
\usepackage{longtable}
\usepackage{calc}
\usepackage{multirow}
\usepackage{hhline}
\usepackage{ifthen}
\usepackage{lscape}

\begin{document}

\bibliographystyle{IEEEtran}

\title{5.5.6}
\author{EE25BTECH11054 - S.Harsha Vardhan Reddy}
\maketitle
\vspace{-1cm}

\renewcommand{\thefigure}{\theenumi}
\renewcommand{\thetable}{\theenumi}
\setlength{\intextsep}{10pt}

\numberwithin{align}{enumi}
\numberwithin{figure}{enumi}
\renewcommand{\thetable}{\theenumi}

\textbf{Question:} \\
For $\vec{A} = \myvec{3 & 1 \\ -1 & 2}$, then $14\vec{A}^{-1}$ is given by

\textbf{Solution:}

Given
\begin{align}
\vec{A} = \myvec{3 & 1 \\ -1 & 2}
\end{align}

Let $\vec{A}^{-1}$ be the inverse of $\vec{A}$. Then
\begin{align}
\vec{A}\vec{A}^{-1} = \vec{I}
\end{align}

Augmented matrix of $\augvec{1}{1}{\vec{A} & \vec{I}}$ is given by
\begin{align}
\augvec{2}{2}{3 & 1 & 1 & 0 \\ -1 & 2 & 0 & 1}
\end{align}

Applying elementary row transformations:
\begin{align}
&\augvec{2}{2}{3 & 1 & 1 & 0 \\ -1 & 2 & 0 & 1} \xrightarrow{R_2 \rightarrow 3R_2 + R_1} \augvec{2}{2}{3 & 1 & 1 & 0 \\ 0 & 7 & 1 & 3}
\end{align}

\begin{align}
&\xrightarrow{R_1 \rightarrow 7R_1 - R_2} \augvec{2}{2}{21 & 0 & 6 & -3 \\ 0 & 7 & 1 & 3}
\end{align}

\begin{align}
&\xrightarrow{R_1 \rightarrow \frac{1}{21}R_1, R_2 \rightarrow \frac{1}{7}R_2} \augvec{2}{2}{1 & 0 & \frac{2}{7} & -\frac{1}{7} \\ 0 & 1 & \frac{1}{7} & \frac{3}{7}}
\end{align}

Hence the inverse of matrix $\vec{A}$ is
\begin{align}
\vec{A}^{-1} = \myvec{\frac{2}{7} & -\frac{1}{7} \\ \frac{1}{7} & \frac{3}{7}}
\end{align}

Therefore, $14\vec{A}^{-1}$ is given by
\begin{align}
14\vec{A}^{-1} = 14 \times \myvec{\frac{2}{7} & -\frac{1}{7} \\ \frac{1}{7} & \frac{3}{7}} = \myvec{4 & -2 \\ 2 & 6}
\end{align}

\end{document}